\section{Definición y cronología de la AGI}

\subsection{Qué es la AGI}

Demis Hassabis define la AGI (Artificial General Intelligence) como \textbf{un sistema de inteligencia artificial capaz de simular la totalidad de las capacidades cognitivas humanas}. A diferencia de la «IA estrecha» actual, la AGI necesita contar con una capacidad de razonamiento general que atraviese distintos dominios: no basta con destacar en una sola tarea, sino que debe poder entretejer el conocimiento proveniente de dominios diferentes en un todo coherente, tal como lo hace el cerebro humano.

\begin{importantbox}{La definición central de la AGI}
AGI = inteligencia artificial general capaz de simular la totalidad de las capacidades cognitivas humanas (full human cognitive capabilities). Las características clave incluyen:
\begin{itemize}
    \item Razonamiento entre dominios e integración del conocimiento
    \item Aprendizaje continuo y actualización de la memoria
    \item Planeación de largo plazo y toma de decisiones en varios pasos
    \item Capacidad de razonamiento general consistente
\end{itemize}
Esto no es un avance en algún benchmark en particular, sino la \textbf{cobertura integral} de las capacidades cognitivas.
\end{importantbox}

\subsection{Cronología: dentro de cinco años}

Hassabis afirma con claridad que existe una \textbf{probabilidad muy alta de que la AGI se logre en los próximos cinco años} («very good chance of it being within the next 5 years»). Vale la pena señalar que este juicio es congruente con la predicción que Shane Legg, cofundador de DeepMind, hizo en 2010, cuando anticipó que la AGI llegaría en 20 años; desde entonces han transcurrido justo unos 16 años.

Los dos motores que impulsan la llegada de la AGI son:
\begin{enumerate}
    \item \textbf{Escalamiento del cómputo (Compute Scaling)}: clústeres de cómputo más grandes y hardware más eficiente
    \item \textbf{Progreso algorítmico (Algorithmic Progress)}: nuevos diseños de arquitectura, métodos de entrenamiento y estrategias de optimización
\end{enumerate}

\begin{knowledgebox}{La visión de Shane Legg}
Shane Legg es uno de los cofundadores de DeepMind y ya en 2010 predijo públicamente que la AGI se lograría en 20 años. En ese momento se consideró una predicción extremadamente optimista, pero hoy Hassabis señala que esa cronología se está cumpliendo. Desde el equipo fundador de tres personas (Hassabis, Legg y Mustafa Suleyman) hasta la Google DeepMind actual, esta convicción de largo plazo sobre la AGI ha sido siempre el motor central del equipo.
\end{knowledgebox}

\subsection{Cuantificar el impacto de la AGI}

Hassabis ofrece una formulación cuantitativa impactante:

\begin{importantbox}{«Diez veces la Revolución Industrial, a diez veces la velocidad»}
\textit{«A veces cuantifico la llegada de la AGI como diez veces la Revolución Industrial, ocurriendo a diez veces la velocidad.» («I sometimes quantify the coming of AGI as 10 times the industrial revolution at 10 times the speed.»)}

La Revolución Industrial tardó más de cien años en que la humanidad asimilara sus efectos secundarios (el trabajo infantil, los dolores de la urbanización, los conflictos de clase, etc.), mientras que la transformación que traerá la AGI se comprimirá para completarse en diez años. Esto significa que la sociedad humana apenas tendrá tiempo de margen para adaptarse.
\end{importantbox}

Antes de la Revolución Industrial, la mortalidad infantil llegaba al 40\%; después de la Revolución Industrial, el progreso tecnológico duplicó la esperanza de vida humana y elevó de manera considerable el nivel de vida. Pero este proceso tomó más de un siglo. La promesa de la AGI es comprimir una transformación de magnitud equivalente (o incluso mayor) dentro del ciclo de vida de una sola generación, lo cual representa a la vez una oportunidad enorme y un reto sin precedentes.

\subsection{Resumen de la sección}

Hassabis define la AGI como una inteligencia general capaz de simular la totalidad de las capacidades cognitivas humanas, con la expectativa de lograrla dentro de cinco años. Cuantifica su impacto con la fórmula «diez veces la Revolución Industrial, a diez veces la velocidad», subrayando que la sociedad humana debe prepararse para esta transformación con una urgencia sin precedentes.

%% ==================== Section 2 ====================
